\section{Small two-power and seven-part characters at coefficient five}
\label{sec:coefficient-five-actual-window}

\paragraph{The proved arithmetic input.}
Use the complete original-coefficient theorem AC-G199,
(F5.1)--(F5.8), on its unchanged original $O_5$.
The new application below deletes precisely its zero rational
noncritical terms from the preceding actual-window remainder.
Its critical partners remain in the one existing AC-G191 aggregate.

\paragraph{The same original window and preceding remainder.}
Retain every original datum, page, prescribed pair, coefficient
completion, lattice, normalization and full local condition of
AC-G198-window. In particular $h\geq2$ and $T\geq1$ are integers,
and at source size $u$ the unchanged parameters are
\[
 \begin{gathered}
 H_u=\mathbb Q(\zeta_{2^{h+2}},\zeta_7,\sqrt{-q},\sqrt{-d},
                    \omega,\sqrt[3]{269}),\quad K=\mathbb Q(\omega),\\
 u^{2+8\varepsilon}\leq2^h\leq2u^{2+8\varepsilon},\quad
 T=\lfloor u^{399}\rfloor,\quad0<\varepsilon<1/256.
 \end{gathered}\tag{F5W.1}
\]
Let $Z_u,J_{u,0},J_{u,\rm mix},V_u,Y_u,C_7(u),U_7(u)$ be exactly
the sets of (CAW.1)--(CAW.4), (RMW.2)--(RMW.3), and (S7W.2).
Their full definitions and every previously established scope are
retained. The preceding remainder is (S7W.5).

Use a primitive order-six character $\xi_7$ modulo seven with
values in the prescribed original $O_5$. Its six powers are all
the modulo-seven characters, with $\xi_7(-1)=-1$.
Put $\mathcal D=\{1,\chi_{-4},\chi_8,\chi_{-8}\}$, and define
\[
 \begin{split}
 \mathcal B_e=\{(\delta,j):{}&\delta\in\mathcal D,\ 0\leq j<6,\\
 &4\nmid e\ \hbox{or }\delta\in\{\chi_{-4},\chi_8\}
           \ \hbox{or }(\delta=1\hbox{ and }2\mid j)\},\\
 C_5(u)={}&\{(5,(\delta\xi_7^j)|_{G_K},e):
                              1\leq e<T,\ (\delta,j)\in\mathcal B_e\},\\
 Z_5(u)={}&\{(5,1,e):1\leq e<T,\ 4\mid e\},\qquad
 U_5(u)=C_5(u)\setminus Z_5(u).
 \end{split}\tag{F5W.2}
\]
All undisplayed cubic and auxiliary finite factors are trivial.
For $p(\delta)\in\{0,1\}$ defined by
$\delta(-1)=(-1)^{p(\delta)}$, let
\[
 b\in\{0,1\},\quad b\equiv p(\delta)+j-e\pmod2,\qquad
 \psi^{\rm nc}_{\delta,j,e}=\delta\xi_7^j\chi_{-3}^{\,b}.
 \tag{F5W.3}
\]
Its sign is $(-1)^e$, as required by the preceding noncritical
rational remainder. The exact input supplied by (F5.3) is
\[
 \mathcal S^{\mathbb Q}_5\bigl((E_5/O_5)
           (\psi^{\rm nc}_{\delta,j,e}c_5^{-e})\bigr)=0
                        \qquad((\delta,j)\in\mathcal B_e),
 \tag{F5W.4}
\]
with the full five-adic local condition, every away-inertia kernel,
the original real-place convention and this unchanged ring.
No original-$K$ strict-zero assertion is assumed.

Define
\[
 \begin{split}
 \mathcal R_{399}^{\rm match,7,5}(u)={}&
 \sum_{(\ell,\eta,e)\in J_{u,\rm mix}}
                         \frac{\log\#\mathcal S_{\ell,\eta}(e)}{b_\ell}\\
 &+\sum_{(\ell,\eta,e)\in
       J_{u,0}\setminus(V_u\cup Y_u\cup C_7(u)\cup C_5(u))}
      \frac{\log\#\mathcal S^{\mathbb Q}_{\ell,
                              \psi_{\eta,e}^{\rm nc}}(e)}{b_\ell}.
 \end{split}\tag{F5W.5}
\]

\begin{theorem}[AC-G200-window: the actual coefficient-five reduction]
\label{AC-G200-window}
All displayed finite characters are actual original-host characters,
and
\[
 \begin{gathered}
 C_5(u)\cap(Z_u\cup V_u\cup Y_u\cup C_7(u))=Z_5(u),\\
 U_5(u)\subset J_{u,0}\setminus(V_u\cup Y_u\cup C_7(u)),\\
 \mathcal R_{399}^{\rm match,7,5}(u)
                      =\mathcal R_{399}^{\rm match,7}(u).
 \end{gathered}\tag{F5W.6}
\]
For $N=T-1$ the exact counts are
\[
 \begin{aligned}
 |C_5(u)|&=24N-9\lfloor N/4\rfloor,\\
 |Z_5(u)|&=\lfloor N/4\rfloor,\\
 |U_5(u)|&=24N-10\lfloor N/4\rfloor.
 \end{aligned}\tag{F5W.7}
\]
Thus the final line is precisely the number of newly absent
rational noncritical terms relative to (S7W.5).
The original triples no longer represented in (F5W.5) form
the disjoint union
$Z_u\sqcup V_u\sqcup Y_u\sqcup U_7(u)\sqcup U_5(u)$.
This is a reduction of the indexing set, not a strict numerical
decrease or a zero assertion for original positive derived indices.

Retain the exact costs and constants of (CAW.5) and (S7W.8).
Then
\[
 \begin{aligned}
 |I_{399}-\mathcal R_{399}^{\rm match,7,5}|&\leq
           \mathcal E_u+\mathcal L_u+\mathcal A_u+C_2\mathcal P_u,\\
 |I_{399}^{\rm all7,2}-\mathcal R_{399}^{\rm match,7,5}|&\leq
           \mathcal E_u+\mathcal L_u+C_2\mathcal P_u.
 \end{aligned}\tag{F5W.8}
\]
There is no new cost or family multiplier. The one full local
majorant, one torsion sum, one critical aggregate and the
$O_\varepsilon(u^{806+16\varepsilon}(\log u)^2)$ envelope remain.
The old and new explicit numerical remainder premises use the
identical window, threshold and slack. Their comparison with the
full original index premise retains the same threshold increase
and half-slack argument.
\end{theorem}

\begin{proof}
\proofstep{1. The actual original coefficient and noncritical lift.}
The two-parts have conductor dividing eight and the seven-parts
have conductor dividing seven. Both cyclotomic factors are in
the same host (F5W.1). Their orders divide $3\cdot2^h$.
The required sixth roots already lie in the original $O_5$,
although the cubic roots need not lie in $\mathbb Z_5$ itself.
There is no new coefficient extension or choice of a different
residue embedding.

Equation (F5W.3) gives sign $(-1)^e$ and hence the literal
noncritical rational lift of the original $K$ coefficient.
If $\tau_5$ denotes the Teichm\"uller lift of $\overline c_5$,
the residual full coefficient is
$\overline\delta\,\overline\xi_7^{\,j}
 \overline\tau_5^{-e}\overline\chi_{-3}^{\,b}$.
The factor $\tau_5$ is used only to identify the residual proof
field; it is not introduced as an original finite host label.
The proved (F5.3), with the full (F5.4) and (F5.8) application,
supplies (F5W.4) on this exact ring and all these local conditions. At five, the original index-two
induced lattice splits because two is a unit in $O_5$.
The inert completion $K_5/\mathbb Q_5$ is retained in the full
local Shapiro comparison; it is not replaced by a split-place
argument. The original critical component remains in the one
G191 aggregate.

\proofstep{2. Every old overlap is an inert-trivial row.}
All $C_5(u)$ triples have coefficient five, trivial cubic part
and trivial auxiliary factors. They therefore miss every old
coefficient-two, coefficient-three and coefficient-$269$ family,
as well as $V_u$ and the coefficient-seven set $C_7(u)$.
Outside those old coefficient families, the only possibilities
in $Z_u$ are the high-two families of (OPW.3) and the inert-trivial
family of (I2W.3). At five the former has threshold
\[
 v_2(5^2-1)=3,\qquad21\cdot5\cdot2^6=6720,
 \qquad c_5=13.
 \tag{F5W.9}
\]
Its primitive two-part conductor is at least $2^{13}$, whereas
the new two-parts have conductor at most eight. These are disjoint.

Five is inert modulo three. The old inert-trivial coefficient-five
family is exactly $(5,1,e)$ with $4\mid e$, namely $Z_5(u)$.
Each such row belongs to $C_5(u)$ because $\delta=1,j=0$
is eligible also when $4\mid e$.
There are no overlaps with $Y_u$: its split-trivial family excludes
five, its residual-matched seven row has coefficient seven, and
its auxiliary coefficients satisfy $q,d>7$.
Distinct two-parts and seven-parts stay distinct over $K$:
their ratio has conductor prime to three and cannot be $\chi_{-3}$.
This proves the first two lines of (F5W.6) exhaustively.

\proofstep{3. Counts and exact deletion.}
At every weight not divisible by four all $4\cdot6=24$ pairs
are eligible. At a weight divisible by four the permitted pairs
are the six $\chi_{-4}$ rows, the six $\chi_8$ rows, and the
three $\delta=1$ rows with $j$ even. Thus precisely nine of the
twenty-four possible pairs are excluded at each such weight.
Among the fifteen permitted rows, exactly one is the old trivial
row $\delta=1,j=0$. This proves (F5W.7), including $T=1$.
For $T\geq2$, the triple $(5,1,1)$ witnesses proper reduction
of the previously retained indexing set.

Apply (F5W.4) to every newly removed rational term. All sums
are finite by the retained G191/G192 finiteness proof, and every
normalization $b_5$ is unchanged. Deleting these zero terms
proves the final equality of (F5W.6), with no additional error.
The nine excluded rows at weights divisible by four are retained.
No original-$K$ strict group or critical partner is made zero
by this composition.

\proofstep{4. One inherited error and the identical common window.}
Substitute the exact equality (F5W.6) into (S7W.8).
This gives (F5W.8) without a second critical aggregate,
second local majorant, second torsion sum or count multiplier.
The original degree $D_u=144\cdot2^h$ and abelian degree
$d_u=96\cdot2^h$ retain their separate roles.

Fix the original datum, page and $0<\varepsilon<1/256$.
Retain exactly the (TC.3) parameters $J,X_0,s$, prescribed pairs
and window $X^{1/J}\leq W$, $W^2\leq X$.
Equal remainders have the identical numerical premise on that
same window, with identical threshold and slack.
For the full index comparison take the same uniform $U_0$ making
the first error of (F5W.8) at most
$(s/2)\exp(u^\varepsilon/2)$, raise the threshold to
$\max(X_0,U_0^J)$, and use that same $W$ with slack $s/2$.
No separate existential windows are intersected. Every original
source rank $R.t$, coefficient, page, pair and simultaneous
finite condition is retained, as is the finite-etale coefficient
and strict completion/inertia dictionary.
\end{proof}

\paragraph{Unchanged open scope.}
The nine unsupported coefficient-five rows at weights divisible
by four, the three unsupported coefficient-seven rows from G198,
and all other mixed and noncritical terms remain. The uniform
numerical remainder bound, mixed theta/Selmer realization and
H04.4 remain open. This composition preserves every preceding
source hypothesis and asks for no new question row. Checks and
actual completion chronology are recorded in the companion events.
